\section{Conclusion}

We identify frequency allocation as a fundamental but under-explored design axis for rotary position embedding and derive EVQ-cosh, a closed-form variational family in which geometric RoPE is the degenerate \(\tau{=}0\) limit. EVQ yields consistent gains across architectures: on MHA, EVQ+YaRN substantially outperforms Geo+YaRN, extending functional context to \(48\mathrm{K}\) (\(24\times\) training length); on MLA, where only 16 frequency channels carry positional information, EVQ alone outperforms Geo+YaRN at \(2\times\) extrapolation; preliminary cross-modal transfer to video DiT provides supporting evidence that the mechanism extends beyond text.

The practical lesson is that long-context systems depend not only on inference-time scaling rules, but crucially on a better training-time frequency substrate. EVQ is a one-line initialization change that adds zero parameters and composes with existing methods, yet on MLA it outperforms the strongest inference-time baseline. As architectures compress RoPE into fewer channels---16 in MLA, potentially fewer in future designs---each channel's placement becomes proportionally more critical, making principled allocation a prerequisite rather than a refinement.
